\section{Ethics, Bias, Privacy, and Reproducibility}
\label{sec:ethics}

\subsection{Data provenance, privacy and licensing}

The three datasets are published for research under MIT licences and consist of textbook diagrams, problem text and annotations; they contain no personal data, so no human-subject or privacy approval applies to the data themselves. The figures originate in commercial textbooks and are used here as the datasets' authors released them, for research; the project redistributes none of the images. Outputs that embed dataset images (the gold-labelling page, cached overlays) are excluded from version control, and the public demonstration of the model shows no dataset image. The synthetic diagrams are generated and carry no rights beyond the project's own. The only human labelling in the study, the gold set, was done by the author on the diagrams alone.

\subsection{Bias in the data and the labels}

The main bias is the one the study is about. Textbook annotation follows textbook convention: a relation is marked when a symbol is printed, and symbols are printed when the problem needs them. A model trained and scored on such labels learns to find marks, and it is penalised, under that score, exactly when it sees a true relation nobody marked (the audit of Section~\ref{sec:intro} puts this at \pctUnmarkedGParallel\% of true parallels in Geometry3K). Scoring against the verifier removes this bias from the evaluation, but it introduces a judgement of its own: the tolerances. The study treats them as a measured choice (calibrated on distributions, validated against a human reading with a pre-registered gate, recorded with the raw labels and every adjudication), and it reports results at the fixed tolerances rather than tuning them to the model. Diagrams that are not drawn to scale are a second, related bias: there, the marks are the truth and the drawing is not; the flag keeps them visible as a category instead of letting them silently lower a score. The synthetic data has a style of its own (clean lines, matplotlib fonts), which could bias a model towards it; the \ftRealFractionPct\%-real batches and the held-out real source are the safeguards, and the extrapolation set is synthetic by design so that a gain there is attributed to size, not style. Finally, the real data comes from English-language, largely United States textbooks, and the vocabulary is six relations; the conclusions are about that domain.

\subsection{Risks in downstream use}

The intended uses are tutoring, automatic grading and accessibility (describing diagrams to blind students). The costs of error differ: a tutor that misreads a diagram explains the wrong problem, a grader penalises a correct observation, a screen reader misinforms a user who cannot check it. Three provisions follow. The model outputs a probability per relation and its calibration is measured (Section~\ref{sec:evaluation}), so that a downstream system can abstain below a confidence it chooses. The verifier is deterministic and inspectable: wherever coordinates are available, a prediction can be checked rather than trusted, and the stored margins say how close a case is to the tolerance. And the model's output is a relation, not a grade or an explanation, so a human remains in the loop for any consequential decision. The model is documented with a model card \cite{mitchell2019modelcards} and the data preparation with a datasheet-style record \cite{gebru2021datasheets} in the repository.

The highlight overlay also makes the question the model is answering visible in the input image; the attention maps of the analysis plan show where it looks; and the asserted and geometric scores are always reported together, so a reader can see how much of the model's ``error'' under one reference is a disagreement between the references.

\subsection{Reproducibility}

The study follows the practices of the reproducibility literature \cite{pineau2021reproducibility} as concrete provisions in the repository. \emph{Code and configuration}: one entry point launches every experiment from a YAML configuration file, with every tunable (paths, backbone, resolution, tolerances, augmentation, optimiser, budget, seeds) in one base file that other configurations override; the full configuration is copied into every run directory. \emph{Determinism}: the seed drives PyTorch, NumPy and Python's random state, the data order, augmentation and head initialisation, with deterministic kernels enabled; three seeds are run for every reported number. \emph{Data}: downloads are verified by SHA-256 and the archive digests are recorded; the processed data, splits, query index and overlays are rebuilt from the raw archives by one command; synthetic diagrams are a pure function of the seed; the unlabelled pool's annotation is sealed with a recorded digest before any experiment. \emph{Runs}: every run writes its metrics, per-epoch history and predictions, and appends to a run log with the configuration hash, seed and best epoch; every table and figure in the reports is regenerated by a script from those files, and every number in this report is produced by a script from the configuration and tables, so that no quantity is typed by hand. \emph{Testing}: every module has a test that runs in under ten seconds, and a smoke configuration exercises the whole pipeline on 50 diagrams before any full run. \emph{Environment}: dependencies are pinned, the hardware and measured throughput are recorded, and the same code runs on the laptop, the laboratory GPUs and Colab with only a path overlay. \emph{Audit trail}: a dated changelog records every decision and every negative finding, and fixed decisions (relation vocabulary, tolerances, splits, budget) may be changed only by editing the plan first and saying so. \emph{Generative AI}: large language model assistance was used for checking references, collecting numbers from the repository and editing prose, and is acknowledged in the submission; the research design and every judgement in this report are the author's own.
